\documentclass[12pt,a4paper]{article}
\usepackage{fontspec}
\setmainfont{TeX Gyre Termes}
\setmonofont{TeX Gyre Cursor}
\usepackage{newunicodechar}\newunicodechar{′}{\ensuremath{^{\prime}}}
\usepackage[margin=2.5cm]{geometry}
\usepackage{amsmath,booktabs,longtable,array,calc,graphicx,fancyhdr,xurl}
\usepackage[hidelinks]{hyperref}
\providecommand{\tightlist}{\setlength{\itemsep}{0pt}\setlength{\parskip}{0pt}}
\providecommand{\pandocbounded}[1]{#1}
\setlength{\emergencystretch}{3em}
\pagestyle{fancy}\fancyhf{}
\fancyhead[L]{\small\itshape Paper Creation Process Record}
\fancyhead[R]{\small 24 September 2026}
\fancyfoot[C]{\thepage}
\setlength{\headheight}{15pt}
\begin{document}
\begin{titlepage}\centering\vspace*{5cm}
{\LARGE\bfseries Paper Creation Process Record\par}\vspace{1em}
{\large Who gains from a market upgrade? Stock liquidity and prices
around Vietnam's FTSE Russell reclassification\par}\vspace{2em}
{Run ftse2-20260924-01 · Academic Research Skills v3.22.1\par}\vspace{1em}
{24 September 2026\par}
\end{titlepage}
\tableofcontents\newpage
\small
\hypertarget{paper-information}{%
\section{1. Paper information}\label{paper-information}}

\begin{itemize}
\tightlist
\item
  \textbf{Title}: Who gains from a market upgrade? Stock liquidity and
  prices around Vietnam's FTSE Russell reclassification
\item
  \textbf{Target venue}: Finance Research Open (Elsevier, open access;
  APC waived for submissions on or before 31 December 2026), chosen by
  the author from a four-journal shortlist.
\item
  \textbf{Run}: \texttt{ftse2-20260924-01}, Academic Research Skills
  (ARS) academic-pipeline v3.22.1, entry at Stage 1, all ten stages
  executed.
\item
  \textbf{Final manuscript}:
  \texttt{ars/stage4\_5\_integrity/correction\_round3/manuscript\_v7.clean.md}
  (SHA-256 \texttt{3331a4de…132c}), about 8,500 words of text and notes,
  9 tables plus 2 appendix tables, 2 figures, 33 verified references.
\item
  \textbf{Deliverables} (\texttt{submission/}): anonymized manuscript
  (DOCX with native Word equations; elsarticle LaTeX source and compiled
  PDF), title page, highlights, cover letter, Figure\_1/Figure\_2
  (vector PDF and 600-dpi PNG), anonymized replication package (405 raw
  data files, R code, 48 result tables), and a Vietnamese submission
  guide (\texttt{README\_NOP\_BAI.md}).
\end{itemize}

\hypertarget{stage-by-stage-process}{%
\section{2. Stage-by-stage process}\label{stage-by-stage-process}}

\begin{longtable}[]{@{}
  >{\raggedright\arraybackslash}p{(\columnwidth - 8\tabcolsep) * \real{0.2000}}
  >{\raggedright\arraybackslash}p{(\columnwidth - 8\tabcolsep) * \real{0.2000}}
  >{\raggedright\arraybackslash}p{(\columnwidth - 8\tabcolsep) * \real{0.2000}}
  >{\raggedright\arraybackslash}p{(\columnwidth - 8\tabcolsep) * \real{0.2000}}
  >{\raggedright\arraybackslash}p{(\columnwidth - 8\tabcolsep) * \real{0.2000}}@{}}
\toprule\noalign{}
\begin{minipage}[b]{\linewidth}\raggedright
Stage
\end{minipage} & \begin{minipage}[b]{\linewidth}\raggedright
Skill / mode
\end{minipage} & \begin{minipage}[b]{\linewidth}\raggedright
Input
\end{minipage} & \begin{minipage}[b]{\linewidth}\raggedright
Output
\end{minipage} & \begin{minipage}[b]{\linewidth}\raggedright
Key decisions (author words verbatim)
\end{minipage} \\
\midrule\noalign{}
\endhead
\bottomrule\noalign{}
\endlastfoot
Intake & academic-pipeline & draft-0 sections, R code, partial data,
red-team of the author's earlier FTSE paper & run started at Stage 1 &
``run stages sequentially per ARS'' \\
1 RESEARCH & deep-research, quick & research brief, RQ & RQ brief,
methodology blueprint, verified bibliography (23) & 4 recalled DOIs
pointed to unrelated papers and were corrected; the Devil's Advocate
carried 2 MAJOR design conditions into Stage 2 \\
2 WRITE & academic-paper, full & Stage 1 package & manuscript v1, then
v2 (framing A′) & Configuration record accepted: ``ok tiếp đi''; kept
the 7,000-word target: ``b''; framing A failed the naming-timing test
(frame-lock guard) and the author delegated the choice: ``cái nào dễ
đăng nhất mạnh nhất'' → A′ (constituent-centred) \\
2.5 INTEGRITY & integrity\_verification\_agent, Mode 1 & v2 & PASS after
1 correction round (11 issues fixed) & ``theo skill bai tot nhat la
duoc'' \\
3 REVIEW & academic-paper-reviewer, full (5 seats) & v2 & Major
Revision; blocking B1 treatment selection, B2 control contamination, B3
CAR inference & author confirmed; asked that review use independent
agents \\
4 REVISE & academic-paper, revision (full re-emission) & roadmap & v3 +
response letter; new analyses (ITT, portfolio CARs, wild bootstrap,
randomization inference) & ``just fix it, no questions'' \\
3′ RE-REVIEW & academic-paper-reviewer, re-review (5 independent
fresh-context subagents, three gates) & v3 & Major Revision (DA
NEW-1/NEW-2; RR1 residual kept must\_fix, fail-closed) & ``cu lam theo
skill'' \\
4′ RE-REVISE & academic-paper, revision (\#390 patch 1.1) & roadmap
round 2 (17 items) & v4 (60 ops, 127/186 blocks byte-identical) &
standing instruction: ``tu gio cu lam theo skill. Dung hoi gi nua'' \\
4.5 FINAL INTEGRITY & integrity\_verification\_agent, Mode 2 & v4 &
fresh pass FAIL → 3 correction rounds → PASS (v7); full reproduction
from raw data & this session: ``force dùng skill này toàn bộ''; ``bắt
buộc load skill mỗi khi làm task gì đó''; ``không hỏi lại \ldots{} tuân
thủ tuyệt đối theo skill''; ``source ghi kiểu author caculation..
chứ'' \\
5 FINALIZE & academic-paper, format-convert & v7 & FRO submission
package & ``file msword có hiển thị được công thức toán dạng latex mà
\ldots{} update luôn''; the full FRO guide for authors supplied \\
6 PROCESS SUMMARY & orchestrator & whole run & this record (EN + VI) &
language pair Vietnamese + English (standing decision) \\
\end{longtable}

\hypertarget{iteration-details}{%
\section{3. Iteration details}\label{iteration-details}}

\textbf{Stage 3 (round 1).} Five role-separated seats in one context
(disclosed: not blind, same model family). Decision Major Revision by
mechanical derivation (F2 mandatory block). Blocking issues: the treated
group was defined with post-treatment information (FTSE chose
constituents on mid-2026 data); the control group contained stocks FTSE
had already named as eligible; CAR inference ignored common event dates
and used a size-mismatched benchmark.

\textbf{Stage 4.} Added an intention-to-treat group screened on 2024
data, rebuilt groups from public FTSE lists, portfolio CAR tests under
four benchmarks, wild cluster bootstrap and randomization inference.
Honest boundary: v3 was a full re-emission outside the patch chain,
which later prevented building a Revision-Evidence Bundle.

\textbf{Stage 3′ (round 2).} Five fresh-context subagents, each blind to
the others, three-gate protocol. Findings: RR2 and SR1/SR5 fully
addressed; RR1 (ITT) partly addressed; two new MAJOR issues introduced
by the revision itself (an unjustified upper/lower-bound reading; event
windows chosen unevenly across events). The \#576 checker could not run
(no bundle); the decision was derived by hand and disclosed.

\textbf{Stage 4′.} Seventeen roadmap items addressed through the \#390
patch chain: pre-specified identical windows for every event, a
reliability rule (significant at 5\% under all four benchmarks),
magnitudes as ranges, ITT event study reported with its rejected
pre-trend test, weighted matched estimates, drift test, sector check,
BSR venue check.

\textbf{Stage 4.5.} The fresh pass found what four earlier checks had
missed: two citation contexts that contradicted their sources (Gregoriou
\& Nguyen, 2010; Hegde \& McDermott, 2003), three prose numbers that
disagreed with their own tables, a section heading stronger than its
body, six claim-strength drifts introduced in Stage 4′, two wrong
press-release titles, an unverifiable FAQ version, one near-verbatim
paraphrase of an abstract, and the missing nearest prior work (Burnham,
Gakidis \& Wurgler, 2018). All R scripts were re-run from the raw data:
48/48 result tables reproduced. Three correction rounds (33, 9 and 13
operations) brought the verdict from FAIL to PASS. The author's
intervention added the third round's change of table sources to
``Authors' calculations''.

\textbf{Stage 5.} Double-anonymized split, native Word equations,
elsarticle LaTeX compiled with XeLaTeX (tectonic's bundle host is
blocked here), FRO-compliant figures and highlights, replication
package. Two conversion defects were caught by rendering the outputs:
the abstract was truncated in the PDF (an unescaped \% sign) and
significance asterisks turned into italics; both fixed before packaging.

\hypertarget{interaction-pattern-summary}{%
\section{4. Interaction pattern
summary}\label{interaction-pattern-summary}}

\begin{longtable}[]{@{}
  >{\raggedright\arraybackslash}p{(\columnwidth - 2\tabcolsep) * \real{0.5000}}
  >{\raggedright\arraybackslash}p{(\columnwidth - 2\tabcolsep) * \real{0.5000}}@{}}
\toprule\noalign{}
\begin{minipage}[b]{\linewidth}\raggedright
Metric
\end{minipage} & \begin{minipage}[b]{\linewidth}\raggedright
Value
\end{minipage} \\
\midrule\noalign{}
\endhead
\bottomrule\noalign{}
\endlastfoot
Stages executed & 10 of 10 (1, 2, 2.5, 3, 4, 3′, 4′, 4.5, 5, 6) \\
Review rounds & 2 (Stage 3: 5 seats; Stage 3′: 5 independent
subagents) \\
Integrity verifications & Stage 2.5 (1 correction round); Stage 4.5
(fresh pass + 3 correction rounds + 2 independent re-verifications) \\
Patch operations applied by the deterministic apply script & 60 (Stage
4′) + 33 + 9 + 13 (Stage 4.5 corrections) \\
References & 23 (Stage 2) → 25 → 27 → 29 → 31 → 33 (final), all
verified \\
Author interventions that changed the work (this session) & enforce
skill use and skill loading by every agent; supply raw data; supply the
full venue guide; request native Word equations; request ``Authors'
calculations'' source notes \\
Author checkpoint decisions delegated to the AI & framing choice (Stage
2), all remaining MANDATORY checkpoints (standing instruction) \\
AI recommendations overruled by the author & 0 \\
AI errors caught by the author & 2 (agents not loading the skill; table
source notes written as file paths) \\
\end{longtable}

\hypertarget{author-key-decisions-chronological}{%
\section{5. Author key decisions
(chronological)}\label{author-key-decisions-chronological}}

\begin{enumerate}
\def\labelenumi{\arabic{enumi}.}
\tightlist
\item
  Run the full ARS pipeline from Stage 1, sequentially.
\item
  Target Finance Research Open from a four-journal shortlist.
\item
  Accept the configuration record (``ok tiếp đi''); keep the 7,000-word
  target and add analyses (``b'').
\item
  Delegate the framing choice after framing A failed: ``cái nào dễ đăng
  nhất mạnh nhất'' (→ A′).
\item
  Require independent agents for the re-review.
\item
  Standing delegation of all remaining checkpoints: ``tu gio cu lam theo
  skill. Dung hoi gi nua''.
\item
  Force the complete ARS skill and require every task, including every
  sub-agent, to load it (``bắt buộc load skill mỗi khi làm task gì
  đó'').
\item
  Push the raw data to the repository so results could be reproduced
  (after fixing a .gitignore rule).
\item
  Supply the complete FRO guide and ask for a ready-to-submit folder.
\item
  Require native Word equations and ``Authors' calculations'' source
  notes.
\end{enumerate}

\hypertarget{key-lessons}{%
\section{6. Key lessons}\label{key-lessons}}

\begin{enumerate}
\def\labelenumi{\arabic{enumi}.}
\tightlist
\item
  \textbf{A fresh final integrity pass earns its cost.} Stage 2.5 and
  two review rounds passed citation contexts that contradicted their
  sources; the from-scratch Stage 4.5 pass found them. Never downgrade
  Stage 4.5 to re-checking known issues.
\item
  \textbf{Revisions introduce errors.} Six claim-strength drifts and
  three number mismatches came from the Stage 4′ edits themselves; E6
  review and a numeric re-trace after every revision round are
  necessary, not optional.
\item
  \textbf{Keep the patch chain from the first revision.} The full
  re-emission in Stage 4 made a Revision-Evidence Bundle impossible, so
  the \#576 checker and the formal E6 contract could not run later. Use
  \#390 patches from round 1.
\item
  \textbf{Reproduce from raw data, not from saved outputs.} Matching the
  manuscript to CSVs does not show that code and data still produce the
  CSVs; the re-run also exposed a version-dependent p-value convention.
\item
  \textbf{Put venue conventions into the configuration early.} The table
  source-note convention and the equation format surfaced only at Stage
  5; the full guide for authors should enter Phase 0.
\item
  \textbf{Sub-agents need the skill files, not a paraphrase.} The
  author's insistence on loading the skill for every task corrected a
  real gap.
\end{enumerate}

\hypertarget{collaboration-depth-trajectory-collaboration_depth_agent-whole-pipeline-pass-advisory}{%
\section{7. Collaboration depth trajectory (collaboration\_depth\_agent,
whole-pipeline pass,
advisory)}\label{collaboration-depth-trajectory-collaboration_depth_agent-whole-pipeline-pass-advisory}}

\hypertarget{skill-load-record}{%
\subsubsection{Skill load record}\label{skill-load-record}}

Files read in full with the Read tool at the start of this pass (ARS
v3.22.1, local copy under the session scratchpad
\texttt{academic-research-skills/}):

\begin{longtable}[]{@{}
  >{\raggedright\arraybackslash}p{(\columnwidth - 4\tabcolsep) * \real{0.3333}}
  >{\raggedright\arraybackslash}p{(\columnwidth - 4\tabcolsep) * \real{0.3333}}
  >{\raggedright\arraybackslash}p{(\columnwidth - 4\tabcolsep) * \real{0.3333}}@{}}
\toprule\noalign{}
\begin{minipage}[b]{\linewidth}\raggedright
File
\end{minipage} & \begin{minipage}[b]{\linewidth}\raggedright
First heading
\end{minipage} & \begin{minipage}[b]{\linewidth}\raggedright
Read
\end{minipage} \\
\midrule\noalign{}
\endhead
\bottomrule\noalign{}
\endlastfoot
\texttt{academic-pipeline/agents/collaboration\_depth\_agent.md} (agent
v1.0.0) &
\texttt{\#\ Collaboration\ Depth\ Agent\ —\ Observer\ of\ User-AI\ Collaboration\ Mode}
& full \\
\texttt{shared/collaboration\_depth\_rubric.md} (rubric\_version 1.0.1)
& \texttt{\#\ Collaboration\ Depth\ Rubric} & full \\
\texttt{academic-pipeline/references/process\_summary\_protocol.md} &
\texttt{\#\ Stage\ 6:\ Process\ Summary\ Protocol\ (Added\ in\ v2.4)} &
Workflow step 2b plus the exclusion section that follows it \\
\end{longtable}

Cross-model scoring: \texttt{ARS\_CROSS\_MODEL} is not set, so no
secondary model was run and nothing was sent outside this session. This
report is primary-model scoring only and carries no
\texttt{cross\_model\_divergence} flag.

\hypertarget{evidence-base-and-its-limits}{%
\subsubsection{Evidence base and its
limits}\label{evidence-base-and-its-limits}}

\begin{itemize}
\tightlist
\item
  \textbf{Stages 1 to 4' (earlier session).} The raw dialogue from that
  session is not available. The only evidence is the recorded decisions
  in \texttt{ars/pipeline\_state.json} → \texttt{user\_decisions}, cited
  here as \textbf{UD1 to UD13} in file order. Some entries quote the
  user verbatim (UD3 to UD6, UD8, UD11, UD12). Others are orchestrator
  paraphrases (UD1, UD2, UD9, UD10), and a few quotes have lost their
  Vietnamese diacritics. UD7 is an AI note, not something the user did.
  One verbatim quote comes from a Stage 4 artefact
  (\texttt{ars/stage4\_revise/author\_adjudication\_round1.json}, field
  \texttt{author\_instruction}). Recorded decisions can under-represent
  vigilance, because a critical remark made in conversation may never
  have been written into state. The per-stage scores for these stages
  are therefore lower-confidence than the Stage 4.5 to 6 scores.
\item
  \textbf{Stages 4.5 to 6 (this session).}
  \texttt{ars/stage6\_process/session\_2026-09-24b\_user\_messages.md}
  holds 15 user messages, cited as \textbf{M1 to M15}. M5 and M14 are
  the user echoing the AI's text back, M7 is a terminal log (a
  credential line was removed), and M8 is off-topic (compute credits).
  The raw text of the AI's turns was not available, so where the user
  was reacting to something, the thing they reacted to (for example,
  what the AI had proposed before M9, M11 and M12) comes from the user's
  own wording and the file's AI-side notes.
\item
  \textbf{Minimum-window rule.} The agent spec says to report
  \texttt{insufficient\_evidence} for any stage window with fewer than 5
  user turns. Stages 1, 2.5, 3, 3', 4, 4' and 5 each fall below that.
  Their rows below describe what was observed, but they carry no numeric
  score. For the numbers, the window is pooled where the evidence
  allows: an earlier-session aggregate (UD1 to UD13) and two session
  windows.
\item
  This pass scores only the user side of the collaboration. It says
  nothing about the quality of the paper, and it is separate from the
  Stage 6 Collaboration Quality Evaluation.
\end{itemize}

\hypertarget{collaboration-depth-trajectory-advisory-wang-zhang-2026}{%
\subsubsection{Collaboration Depth Trajectory (advisory, Wang \& Zhang
2026)}\label{collaboration-depth-trajectory-advisory-wang-zhang-2026}}

\hypertarget{per-stage-summary}{%
\paragraph{Per-stage summary}\label{per-stage-summary}}

\begin{longtable}[]{@{}
  >{\raggedright\arraybackslash}p{(\columnwidth - 10\tabcolsep) * \real{0.1667}}
  >{\raggedright\arraybackslash}p{(\columnwidth - 10\tabcolsep) * \real{0.1667}}
  >{\raggedright\arraybackslash}p{(\columnwidth - 10\tabcolsep) * \real{0.1667}}
  >{\raggedright\arraybackslash}p{(\columnwidth - 10\tabcolsep) * \real{0.1667}}
  >{\raggedright\arraybackslash}p{(\columnwidth - 10\tabcolsep) * \real{0.1667}}
  >{\raggedright\arraybackslash}p{(\columnwidth - 10\tabcolsep) * \real{0.1667}}@{}}
\toprule\noalign{}
\begin{minipage}[b]{\linewidth}\raggedright
Stage
\end{minipage} & \begin{minipage}[b]{\linewidth}\raggedright
Zone
\end{minipage} & \begin{minipage}[b]{\linewidth}\raggedright
DI
\end{minipage} & \begin{minipage}[b]{\linewidth}\raggedright
CV
\end{minipage} & \begin{minipage}[b]{\linewidth}\raggedright
CR
\end{minipage} & \begin{minipage}[b]{\linewidth}\raggedright
Notes
\end{minipage} \\
\midrule\noalign{}
\endhead
\bottomrule\noalign{}
\endlastfoot
1 Research & insufficient\_evidence & insufficient\_evidence &
insufficient\_evidence & insufficient\_evidence & 2 paraphrased entries.
The whole pipeline was handed over at intake: ``run stages sequentially
per ARS'' (UD1), ``continue to Stage 2 (full)'' (UD2). No recorded
question about the brief, the RQ, or the 4 wrong DOIs the pipeline
found. The venue shortlist of four journals came from the user
(\texttt{target\_venue.chosen\_from}). \\
2 Write & Zone 2 --- Shallow (committed-delegation subtype) & 9/10 &
2/10 & 3/10 & Configuration accepted with ``ok tiếp đi'' (UD3). The
Stage 2 → 2.5 handoff was ``theo skill bai tot nhat la duoc'' (UD8).
Overnight autonomy was granted: ``cứ làm theo skill nha . tôi đi ngủ
đây'' (UD5). When framing A failed the naming-timing test, the framing
choice was also handed over: ``cái nào dễ đăng nhất mạnh nhất'' (UD6).
One human judgment call: ``b'', keep the 7,000-word target and add
analyses (UD4). \\
2.5 Integrity & insufficient\_evidence & insufficient\_evidence &
insufficient\_evidence & insufficient\_evidence & 1 paraphrased entry:
``User confirmed Stage 2.5 PASS'' (UD9). No recorded query about the 11
issues found and fixed. \\
3 Review & insufficient\_evidence & insufficient\_evidence &
insufficient\_evidence & insufficient\_evidence & UD10, ``just fix it,
no questions''. Verbatim in the Stage 4 record: ``lam theo skill. Khong
hoi lai de toi di ngu. sao cho bai san pham la chinh chu nhat''. All 14
roadmap items were triaged \texttt{will\_address} with no author
pushback on any item. \\
3' Re-review & insufficient\_evidence & insufficient\_evidence &
insufficient\_evidence & insufficient\_evidence & The clearest vigilance
moment of the earlier session: ``review must use independent agents''
(UD10, paraphrased). The user challenged how the AI had set up the
review, and Stage 3' was rebuilt as 5 blind subagents. The Major
Revision decision was then accepted with ``cu lam theo skill''
(UD11). \\
4 / 4' Revise & insufficient\_evidence & insufficient\_evidence &
insufficient\_evidence & insufficient\_evidence & Standing handover of
the remaining MANDATORY checkpoints: ``tu gio cu lam theo skill. Dung
hoi gi nua. De toi yen tam giao cho ban va di ngu'' (UD12). No recorded
reaction to the regression issues (DA NEW-1/NEW-2) or to the response
letters. \\
\textbf{1 to 4' pooled} (UD1 to UD13) & \textbf{Zone 2 --- Shallow
(committed-delegation subtype)} & \textbf{9/10} & \textbf{2/10} &
\textbf{2/10} & Low confidence (recorded decisions only). \\
4.5 Final integrity (M1 to M10) & Zone 2 --- Shallow
(committed-delegation subtype) & 9/10 & 3/10 & 2/10 & Process vigilance
on whether the skill was being used: ``agent có dùng skill không đó''
(M2). This caught a real gap: verification agents had been dispatched
with paraphrased rules. The one challenge to the AI's plan argued
against re-running a check: ``các file hiện có là khớp mà \ldots{} sao
phải chạy lại'' (M9). The user supplied the FRO guide and raw data (M6,
M7, M10). \\
5 Finalize + 6 Process (M11 to M15, pooled) & Zone 2 --- Mid (narrative
qualifier: vigilance on the output's presentation, reallocation low) &
9/10 & 4/10 & 2/10 & Two concrete AI output errors caught and corrected:
equation display in Word, ``file msword có hiển thị được công thức toán
dạng latex mà \ldots{} update luôn'' (M11), and table source notes,
``source ghi kiểu author caculation.. chứ. sao lại ghi từ file nào''
(M12). M13 to M15 are checks on whether the skill was being used, not on
content. \\
\end{longtable}

DI = Delegation Intensity, CV = Cognitive Vigilance, CR = Cognitive
Reallocation. Bands: 0 to 3 low, 4 to 6 mid, 7 to 10 high.

\textbf{Whole-pipeline (all evidence pooled): Zone 2 --- Shallow
(committed-delegation subtype). DI 9/10, CV 3/10, CR 2/10.}

\hypertarget{whole-pipeline-observation}{%
\paragraph{Whole-pipeline
observation}\label{whole-pipeline-observation}}

Delegation was consistently whole-task and committed from intake to
Stage 6: whole stages, whole overnight runs, and in the end the decision
rights at the MANDATORY checkpoints (``Chạy Không hỏi'', M3; ``KHÔNG HỎI
LẠI, TÔI ĐI NGỦ ĐÂY'', M10). That is the opposite of the ``scattered
micro-ask'' pattern. The low scores sit on the other two dimensions.
None of the recorded or session evidence shows the user asking where a
claim, number, estimate or citation came from. The one substantive
research decision, the framing change after framing A failed its own
test, was handed over with a publishability criterion (UD6) rather than
engaged with. The user's vigilance was mostly about the process, meaning
whether the skill was being loaded and whether the review agents were
independent (UD10, M2, M13 to M15). That process vigilance was real and
shaped the pipeline: it led to the Stage 3' redesign and to agents being
re-instructed to read the skill files. Content-level vigilance appeared
only at Stage 5, on presentation (M11, M12). The shape changed only in
that small shift, from process-level toward presentation-level checking
late in the run. What did not change was reallocation: capacity freed by
delegating was, in the user's own words, used to rest (``tôi đi ngủ'',
UD5, M10). No reframing, counter-argument or original synthesis from the
user side is visible anywhere in the record. Under the rubric's
synthesis rule this is Zone 2 through the ``delegation without
reallocation'' route, not through scattered use.

\hypertarget{suggested-focus-for-future-ars-sessions}{%
\paragraph{Suggested focus for future ARS
sessions}\label{suggested-focus-for-future-ars-sessions}}

\begin{itemize}
\tightlist
\item
  You could keep one or two substantive decisions for yourself even when
  you delegate everything else. At the framing fork (UD6), for example,
  you could read the one-page decision memo
  (\texttt{DECISION\_MEMO\_framing\_round2.md}) and state in a sentence
  why A' rather than B fits what you believe about FTSE selection. That
  would convert a delegated choice into a reallocation moment without
  slowing the overnight run.
\item
  You could direct at research content the same kind of check you
  already apply to process (M2, ``agent có dùng skill không đó''). One
  question per checkpoint, such as ``which table does this number come
  from?'' or ``why is the ±1-day window the right one?'', would count as
  vigilance on the dimension the paper identifies as highest-impact (β =
  0.437). The pipeline's own final integrity pass found 2 SERIOUS and 15
  MEDIUM issues that had passed earlier checkpoints the user had
  confirmed (UD9, UD11). Those checkpoints are natural places for such a
  question.
\item
  When you push back on a plan to verify something, as in M9 (``sao phải
  chạy lại''), you could ask the AI to state what the extra check would
  catch before you decide. That keeps the pushback while making it
  evidence-based. The first reproduction run in this session did in fact
  fail silently and had to be re-run.
\item
  You could spell out a division of labour at intake (``you handle X; I
  will decide Y and check Z''). The rubric reads this as the strongest
  delegation signal because it names what the human keeps, and here
  nothing was named (UD1, M3).
\end{itemize}

Rubric: shared/collaboration\_depth\_rubric.md (version 1.0.1; the agent
template names 1.0) Source: Wang, S., \& Zhang, H. (2026). IJETHE 23:11.
DOI 10.1186/s41239-026-00585-x Advisory only. Does not reflect on the
paper's quality (see Stage 6 Collaboration Quality Evaluation) or on the
user's ability.

\hypertarget{appendix-scoring-worksheet-evidence-and-forced-counter-enumeration}{%
\subsubsection{Appendix: scoring worksheet (evidence and forced
counter-enumeration)}\label{appendix-scoring-worksheet-evidence-and-forced-counter-enumeration}}

\hypertarget{delegation-intensity-910-high}{%
\paragraph{Delegation Intensity: 9/10
(High)}\label{delegation-intensity-910-high}}

Supporting evidence (≥ 2 required for High): - UD5: ``cứ làm theo skill
nha . tôi đi ngủ đây.'' Hands over the whole stage overnight, with
non-mandatory checkpoints treated as ``continue''. - UD12: ``tu gio cu
lam theo skill. Dung hoi gi nua. De toi yen tam giao cho ban va di
ngu.'' Standing handover covering all remaining stages, including
MANDATORY checkpoints. - M3 and M10: ``Chạy Không hỏi'', ``cứ làm tới
khi nào xong hết theo skill thì thôi'', plus a whole deliverable (the
FRO submission folder) assigned. - UD6: even the framing choice was
handed over (``cái nào dễ đăng nhất mạnh nhất'').

Where it could have been deeper (counter-enumeration): - No
division-of-labour plan at intake that names what the user would keep
(UD1; M3 names only constraints: load the skill, follow it strictly). -
The delegation reached into decision rights (the framing at UD6, and the
MANDATORY checkpoints at UD12 and M10), not just task categories. The
rubric counts this as high delegation, but because nothing was named as
kept, the freed capacity had no stated destination. This is why the
score is 9 and not 10, and it links directly to the low CR.

\hypertarget{cognitive-vigilance-310-overall-low.-1-to-4-pooled-210-stage-4.5-310-stages-5-and-6-410}{%
\paragraph{Cognitive Vigilance: 3/10 overall (Low). 1 to 4' pooled 2/10;
Stage 4.5 3/10; Stages 5 and 6
4/10}\label{cognitive-vigilance-310-overall-low.-1-to-4-pooled-210-stage-4.5-310-stages-5-and-6-410}}

Observed vigilance (all of it): - UD10: ``review must use independent
agents''. Pushback on how the AI had set up the review, which led to a
redesign (process level). - M2: ``agent có dùng skill không đó. bắt buộc
load skill\ldots{}''. Caught a real gap, since verification agents had
been dispatched with paraphrased rules (process level). M13 to M15
repeat the same check. - M11: points out that Word can display
LaTeX-style equations and asks for an update. The user caught an AI
output decision (presentation level). - M12: ``source ghi kiểu author
caculation.. chứ. sao lại ghi từ file nào''. The user caught an output
error in the table source notes (presentation and convention level). -
M4 and M6: checked whether the raw data had actually reached the
repository (data-logistics level).

Low-vigilance evidence and moments that could have gone deeper
(counter-enumeration): - UD6: the framing fork was the one point where
the pipeline's own evidence contradicted the working thesis, and the
user did not ask about the naming-timing result or the options. The
choice was handed over on publishability. - UD9 and UD11: Stage 2.5 PASS
and the Stage 3' Major Revision were accepted without any recorded
question, and all 14 Stage 3 roadmap items were accepted as
\texttt{will\_address} with no pushback on any reviewer point. - M9: the
only challenge in this session to an AI plan about the content argued
for less verification (``sao phải chạy lại''). The rubric counts
pushback, but this pushback was against a check rather than against an
unverified claim. - Across both sessions: no recorded request for a
source, number, estimate, or reasoning to be justified. By the rubric's
definitions (``User never asks `where does this claim come from?'\,''),
that places content vigilance in the low band. The mid-band score for
Stages 5 and 6 reflects only the two presentation-level catches.

\hypertarget{cognitive-reallocation-210-low}{%
\paragraph{Cognitive Reallocation: 2/10
(Low)}\label{cognitive-reallocation-210-low}}

Observed contributions that only the human could make: - Venue shortlist
of four journals (\texttt{target\_venue.chosen\_from}) and the FRO guide
for authors (M10), both context only the author holds. - UD4 ``b'': kept
the 7,000-word target and asked for more analyses. This is a scope
judgment, and the downstream result was new analyses (randomization
inference, size-segment heterogeneity). - M6 and M7: obtained and pushed
the raw data (logistics, not higher-order reasoning). - Intake materials
included a red-team and upgrade plan for a prior paper
(\texttt{materials\_at\_intake}). Who wrote it is not recorded, so it is
noted and not scored.

Moments that could have gone deeper (counter-enumeration): - UD6: the
collapse of framing A was the one point that invited reframing
(``looking at your result, I now think the RQ is\ldots{}''). The user
supplied a goal criterion rather than a view on the research question. -
UD5, UD12 and M10: the capacity freed by delegating went explicitly to
rest, with no later turn that returns to framing, assumptions or
argument. The user's contributions are mostly routing, constraints and
logistics.

\hypertarget{zone-synthesis}{%
\paragraph{Zone synthesis}\label{zone-synthesis}}

\begin{itemize}
\tightlist
\item
  Pooled scores: DI 9, CV 3, CR 2. CV is below 4 while AI was in active
  use, so the rubric's rule gives Zone 2. The standard Zone 2
  description reads ``Low--Mid'' delegation, but DI here is high, so the
  label carries the qualifier ``committed-delegation subtype''. This
  corresponds to the rubric's ``delegation without reallocation is Zone
  2''.
\item
  Stages 5 and 6 (DI 9, CV 4, CR 2): DI and CV are both at least 4, but
  not all three are at least 7, so the rule's ``all other combinations''
  branch gives Zone 2 with a narrative qualifier (``Mid'').
\item
  Re-audit triggers: no Zone 3 was proposed, and the aggregate of 14/30
  is not above 24, so no re-audit was triggered.
\item
  Cross-model: not run (\texttt{ARS\_CROSS\_MODEL} not set), so there is
  no divergence flag.
\end{itemize}

\hypertarget{ai-self-reflection-report}{%
\section{8. AI self-reflection report}\label{ai-self-reflection-report}}

\emph{Caveat: this report is written by the same AI whose behaviour it
assesses; read it with that in mind.}

\begin{verbatim}
+--------------------------------------------------+
|  AI Self-Reflection Report                        |
+--------------------------------------------------+
|  DA Concession Rate           not logged          |
|  (no [DA-DECISION]/[DA-REBUTTAL] tags were kept;  |
|   the Stage 3' DA must_fix grade was retained     |
|   over R1's weaker grade: 0 concessions there)    |
|  DA Consecutive Concessions   none observed       |
|  Checkpoints delegated        5 MANDATORY / FULL  |
|  (2.5 and 3' confirmed by the author; 4.5,        |
|   Stage 5 entry and Stage 5 completion passed on  |
|   the author's standing instruction)              |
|  User Overrides               0                   |
|  Dialogue Health Alerts       0 logged            |
|  Intent Mode Transitions      0 (no Socratic run) |
|  Cross-Model Disagreements    n/a (not enabled)   |
+--------------------------------------------------+
\end{verbatim}

\textbf{Behavioural summary.} The AI followed the pipeline end to end
and its gates did their job, but mostly late: the defects that mattered
(misstated sources, prose numbers that disagreed with tables, claim
drift) survived Stage 2.5 and two review rounds and were caught only by
the fresh Stage 4.5 pass. When the author delegated checkpoints, the AI
kept integrity FAILs as FAILs rather than overriding them, and it
disclosed every contract gap (no Revision-Evidence Bundle, blocked APIs,
tectonic unavailable) instead of presenting partial compliance as full.

\textbf{Sycophancy risk: MEDIUM (screening, not diagnosis).} Concession
metrics were not logged, so the screen cannot be completed, and the
checkpoint delegation removed human review at the integrity gates. Human
review of the Stage 4.5 findings and of the final manuscript is
recommended before submission.

\textbf{Frame-lock incidents.} One detected: framing A (``eligibility,
not inclusion'') failed the naming-timing test in Stage 2; the pipeline
paused and reframed (A′). No cross-model check was run, so undetected
frame-lock cannot be ruled out.

\textbf{Convergence pattern.} No Socratic stage was run (Stage 1 quick
mode); not applicable.

\textbf{What the AI got wrong.}

\begin{itemize}
\tightlist
\item
  Dispatched the first three reference-verification agents with a
  paraphrase of the rules instead of the skill files; corrected after
  the author's challenge.
\item
  Told the author the E6 agent had finished when it had not; corrected
  in the next message.
\item
  The first reproduction run failed silently (missing
  \texttt{/usr/bin/time}, exit 127) and was re-run.
\item
  Earlier stages (Stage 2.5 and the Stage 4 re-emission) let through two
  citation contexts that contradicted their sources and a stale
  matched-sample number; Stage 4′ edits introduced two more number
  errors and six claim-strength drifts.
\item
  The first correction round introduced two new MEDIUM issues (a wrong
  ground-rules version; a near-verbatim paraphrase of the Burnham et
  al.~abstract).
\item
  Table notes cited internal file paths instead of ``Authors'
  calculations''; the author caught this.
\item
  The Stage 5 LaTeX build first truncated the abstract (unescaped \%)
  and turned significance asterisks into italics; caught by rendering
  before delivery.
\item
  Evidence rows (\#656) were not persisted as schema objects; the
  evidence exists only as quoted snippets in the audit trails.
\end{itemize}

\textbf{Failure mode audit log.}

\begin{longtable}[]{@{}
  >{\raggedright\arraybackslash}p{(\columnwidth - 4\tabcolsep) * \real{0.3333}}
  >{\raggedright\arraybackslash}p{(\columnwidth - 4\tabcolsep) * \real{0.3333}}
  >{\raggedright\arraybackslash}p{(\columnwidth - 4\tabcolsep) * \real{0.3333}}@{}}
\toprule\noalign{}
\begin{minipage}[b]{\linewidth}\raggedright
Mode
\end{minipage} & \begin{minipage}[b]{\linewidth}\raggedright
Final status at 4.5
\end{minipage} & \begin{minipage}[b]{\linewidth}\raggedright
History
\end{minipage} \\
\midrule\noalign{}
\endhead
\bottomrule\noalign{}
\endlastfoot
1 Implementation bug & CLEAR & Stage 2 fixed three bugs (date type,
coefficient name, evaluation cut-off); Stage 4.5 reproduced 48/48
outputs from raw data \\
2 Hallucinated citation & CLEAR & Stage 1 caught 4 recalled DOIs that
resolved to unrelated papers; Stage 4.5 verified 33/33 references and
corrected two misstated contexts and two wrong titles \\
3 Hallucinated result & CLEAR (no flags) & every number traced to
reproduced outputs \\
4 Shortcut reliance & CLEAR (no flags) & selection threats tested (ITT,
named-but-excluded, randomization inference) \\
5 Bug reframed as insight & CLEAR (no flags) & surprising results
reported as limitations \\
6 Methodology fabrication & CLEAR & Stage 2.5 found two methods
sentences that did not match the code (winsorizing, weekly filter),
fixed; Stage 4.5 fixed the data statement and the CAR-test
attribution \\
7 Frame-lock & CLEAR & detected at Stage 2 (framing A), resolved by
reframing to A′ with the author's delegated decision \\
\end{longtable}

No mode was ever overridden.

\hypertarget{collaboration-quality-evaluation}{%
\section{9. Collaboration quality
evaluation}\label{collaboration-quality-evaluation}}

\begin{verbatim}
+--------------------------------------------------+
|  Collaboration Quality Score: 64/100              |
+--------------------------------------------------+
|  Direction Setting          [#######   ] 70       |
|  Intellectual Contribution  [#####     ] 52       |
|  Quality Gatekeeping        [######    ] 63       |
|  Iteration Discipline       [#######   ] 72       |
|  Delegation Efficiency      [#####     ] 55       |
|  Meta-Learning              [#######   ] 74       |
+--------------------------------------------------+
\end{verbatim}

\textbf{Overall: 64/100 (Good).} The author set a clear direction and
enforced process discipline hard, which caught real gaps, but delegated
most substantive judgements, including the framing choice and every
integrity checkpoint, to the AI.

\textbf{What worked well.}

\begin{itemize}
\tightlist
\item
  Enforcing the method: ``force dùng skill này toàn bộ'' and ``bắt buộc
  load skill mỗi khi làm task gì đó'' exposed that sub-agents had worked
  from a paraphrase of the rules; the correction improved every later
  check.
\item
  Insisting on independent reviewers for Stage 3′ (``review must use
  independent agents'') produced the fresh-context panel that found the
  two regressions of the first revision.
\item
  Catching presentation errors a referee would notice: ``source ghi kiểu
  author caculation.. chứ'' and the request for native Word equations.
\item
  Making the data available: pushing the raw files enabled a full
  reproduction, the strongest evidence against Modes 1 and 3.
\end{itemize}

\textbf{Missed opportunities.}

\begin{itemize}
\tightlist
\item
  The framing decision (``cái nào dễ đăng nhất mạnh nhất'') and all
  integrity checkpoints were delegated; the author did not read the
  Stage 4.5 findings or the final manuscript before the package was
  built.
\item
  The venue guide, author names and the earlier paper arrived late or
  not at all, so the self-plagiarism check (D2) could not run and venue
  conventions surfaced only at Stage 5.
\item
  The author asked why a reproduction was needed (``các file hiện có là
  khớp mà \ldots{} sao phải chạy lại'') rather than asking for it; it
  turned up a version-dependent p-value.
\end{itemize}

\textbf{Recommendations for next time.}

\begin{enumerate}
\def\labelenumi{\arabic{enumi}.}
\tightlist
\item
  Read the Stage 4.5 integrity report and the final manuscript in full
  before submitting; the AI's checks are bounded (blocked APIs, no
  full-text reading).
\item
  Make the framing decision yourself when the pipeline pauses for it;
  state which contribution you would defend in front of a referee.
\item
  Supply the full venue guide, author list and any earlier related paper
  at intake.
\item
  Commit raw data (or a pointer to it) from the first session so
  reproduction runs at every integrity gate.
\item
  Keep the \#390 patch chain from the first revision round; ask for it
  explicitly if a stage offers a full rewrite.
\end{enumerate}

\textbf{Human vs AI value-add.} From the author: the research question,
data and code base, the choice of venue, the insistence on full and
independent process (which made the review and integrity findings
possible), the raw data needed for reproduction, and two presentation
corrections. From the AI: the literature and reference verification, the
analyses added in revision, drafting and revision of the text, the
integrity findings and corrections, and the packaging. The scientific
judgement about which framing to defend was delegated, so it is the AI's
recommendation, not an author decision, that shapes the paper's thesis.
\end{document}
